%% =================================================================
%% verzeichnisse/abkuerzungen.tex
%% =================================================================

% Nutzungshinweise:
% \gls{tag}     -> Erstes Auftreten: Vollform (Kurzform)
% \glspl{tag}   -> Pluralform
% \acrfull{tag} -> Erzwingt Vollform

% =========================================================
% KI / Sprachverarbeitung / Frameworks
% =========================================================
\newacronym{stt}{STT}{Speech-to-Text}
\newacronym{nlp}{NLP}{Natural Language Processing}
\newacronym{llm}{LLM}{Large Language Model}
\newacronym{api}{API}{Application Programming Interface}
\newacronym{rest}{REST}{Representational State Transfer}
\newacronym{json}{JSON}{JavaScript Object Notation}
\newacronym{sdk}{SDK}{Software Development Kit}

% =========================================================
% Datenverarbeitung & Backend
% =========================================================
\newacronym{sql}{SQL}{Structured Query Language}
\newacronym{sqlite}{SQLite}{Serverlose, eingebettete SQL-Datenbank-Engine}
\newacronym{etl}{ETL}{Extract, Transform, Load}
\newacronym{pipeline}{Pipeline}{Modulare Verarbeitungskette zur schrittweisen Datenverarbeitung}

% =========================================================
% KI-Modelle & Verarbeitung
% =========================================================
\newacronym{whisper}{Whisper}{Automatisches Speech-to-Text-Modell von OpenAI}
\newacronym{spacy}{spaCy}{Bibliothek zur industriellen Verarbeitung natürlicher Sprache}
\newacronym{token}{Token}{Einzelnes linguistisches Element innerhalb eines Textes}

% =========================================================
% UI / UX / Frontend (Flutter / Streamlit)
% =========================================================
\newacronym{ui}{UI}{User Interface}
\newacronym{ux}{UX}{User Experience}
\newacronym{gui}{GUI}{Graphical User Interface}
\newacronym{flutter}{Flutter}{Cross-Platform UI Framework von Google}
\newacronym{widget}{Widget}{Grundbaustein der Benutzeroberfläche in Flutter}

% =========================================================
% Mobile / System
% =========================================================
\newacronym{apk}{APK}{Android Application Package}
\newacronym{os}{OS}{Operating System}
\newacronym{ram}{RAM}{Random Access Memory}
\newacronym{cpu}{CPU}{Central Processing Unit}
\newacronym{gpu}{GPU}{Graphics Processing Unit}
\newacronym{wlan}{WLAN}{Wireless Local Area Network}

% =========================================================
% Audio / Datenquellen
% =========================================================
\newacronym{mic}{MIC}{Mikrofon als Audioeingabegerät}
\newacronym{audio}{Audio}{Digitale Schallsignaldaten als Eingabequelle für STT-Systeme}